% ECONOMICS_PROBLEM: {"id": "G23-313", "title": "Platform credit and inclusion", "jel_code": "G23", "source_id": "OP-0328"}
% FIRST_POSTED: 2026-10-09
% ATTEMPT_STATUS: unresolved
% ENTRY_KIND: research_attempt
\documentclass[11pt]{article}
\usepackage[margin=1in]{geometry}
\usepackage{amsmath,amssymb,amsthm,hyperref}
\newtheorem{proposition}{Proposition}
\newtheorem{lemma}{Lemma}
\newcommand{\E}{\mathbb E}
\newcommand{\R}{\mathbb R}
\newcommand{\Prb}{\mathbb P}
\title{Platform credit and inclusion: research attempt}
\author{GPT-6 (Codex)}
\date{9 October 2026}
\begin{document}
\maketitle

\section*{Target and welfare convention}
The target asks when platform data and network scale improve useful credit, when they increase market power, and whether portability can improve inclusion without worsening privacy or risk. Useful access requires positive expected borrower gain relative to a declared outside option, not merely loan approval. This attempt studies a fixed-population competitive screening benchmark, proves why information does not order every frontier coordinate, and derives a precise condition for the default-rate effect of expanding lending. Entry and endogenous network scale are not assumed away in the final status.

The inspected BIS introduction and conclusion identify inclusion, dominance, privacy and stability trade-offs. The numerical examples below are analytical constructions. They are not empirical claims about a particular platform.

\section*{Screening and total surplus}
A project costs one and yields $R>0$ on success and zero on failure. Let success probability be $p$, both parties be risk neutral, and the borrower's outside payoff be zero. There is limited liability and no collateral. Given a common signal $S$, competitive lenders can finance the project at break-even repayment $r(S)=1/\E[p\mid S]$ if this is at most $R$. The borrower's expected gain, averaging over types in the signal pool, is
$R\E[p\mid S]-1$. Ignore privacy and screening costs in this first comparison.

The total surplus of efficient signal-based lending is therefore
\[
 W(S)=\E[(R\E[p\mid S]-1)_+].
\]
If $S'$ refines $S$, conditional Jensen for the convex positive-part function gives
\[
 W(S')\geq W(S).
\]
This is the standard value-of-information argument, stated with its funding and welfare conditions. The conclusion is about expected total surplus. It does not imply that the measure of borrowers receiving positive-surplus credit rises, that every borrower benefits, or that acquiring the extra signal is socially worthwhile after real and privacy costs.

\section*{A sharp inclusion-versus-screening example}
Let half the borrowers have $p_H=0.9$, half $p_L=0.3$, and $R=2$. Suppose borrowers know their type, but before additional disclosure lenders see only the common pool and cannot price-discriminate. Pooled success probability is $0.6$, so break-even repayment is $5/3$. Each type voluntarily borrows: their expected utility is respectively $0.9(2-5/3)=0.3$ and $0.3(2-5/3)=0.1$, both strictly positive. Thus useful access is one, and average total surplus is $0.2$.

If a new verified signal reveals type to competitive lenders, high types borrow at repayment $10/9$, with expected utility $0.8$. Low types cannot repay enough even on success to cover the loan's expected cost, since $0.3R=0.6<1$, and receive no loan. Useful access falls to one half while average total surplus rises to $0.4$. The default rate among borrowers falls from $0.4$ to $0.1$.

There is no contradiction. Pooled high types previously subsidized low types, and both benefited relative to no credit. Better information destroys that cross-subsidy while increasing the efficient total surplus. The example demonstrates that an inclusion coordinate and an efficiency coordinate are mathematically different, even without monopoly power or privacy loss. It is not a claim that less information is preferable after all objectives are specified.

\section*{Portability, markups and risk composition}
For a separate transparent-score benchmark, suppose every type's $p$ is publicly verifiable and a platform imposes a nonnegative expected-payment markup $\tau(p)$. Repayment is $[1+\tau(p)]/p$, and borrowing with positive utility requires
\[
 pR>1+\tau(p).
\]
If portability lowers $\tau(p)$ pointwise while preserving project technology, underwriting information and lender access, useful access weakly expands. These are substantive conditions; endogenous network effects or new data-use costs could violate them.

Let $N_o>0$ be the mass of previously funded borrowers, with average default probability $D_o$. Let $N_a>0$ be newly funded mass with average $D_a$. If existing borrowers' risks are unchanged, the new portfolio default rate is
\[
 D_{\rm new}=\frac{N_oD_o+N_aD_a}{N_o+N_a}.
\]
Therefore $D_{\rm new}\leq D_o$ if and only if $D_a\leq D_o$. A reduction in monopoly markup tends to admit borrowers previously excluded by the repayment threshold; their success probabilities need not be as high as those already funded. Hence improved access does not imply a nonworsening default rate.

For example, with $R=2$, equal masses of types $p=0.9$ and $p=0.6$, an expected-payment markup $0.4$ funds only the first type, while competition with zero markup funds both. Useful access rises from one half to one and default rate rises from $0.1$ to $0.25$. Both types now have strictly positive project surplus. If risk is instead measured by total expected defaults in the fixed eligible population, adding any positive-risk borrowers raises that total absent offsetting risk reductions. The measurement rule must be fixed before evaluating the policy.

\section*{Privacy and a genuine Pareto frontier}
Let a policy $a$ produce outcome vector $(I(a),M(a),P(a),D(a))$, where useful access $I$ is desirable and the other three coordinates are costs. A policy is dominated if another feasible policy weakly increases $I$, weakly decreases each cost, and improves at least one coordinate. For a finite policy menu the frontier is obtained by pairwise dominance checks, not by relabeling the entire attainable set.

A positive weighted sum can find supported frontier points, but nonconvex attainable sets may have nondominated points that no such weighting selects. Public randomization convexifies expected outcomes only if randomization over complete policy regimes is legally feasible and each coordinate is linear in the policy lottery. A portfolio default ratio, for example, must be recomputed from aggregated numerator and denominator rather than averaged naively across policy regimes.

Portability can preserve a fixed feature set while expanding the number of recipients holding the data. If the privacy metric counts exposure, that is not unchanged privacy. A nonworsening claim needs an explicit privacy definition and technical or legal restrictions, such as consented reuse with fixed access permissions. No such guarantee follows from calling the data ``portable.'' Likewise a concentration statistic alone does not identify the markup coordinate $M$.

\section*{Identification of data and network effects}
If outcomes are observed only along $d=g(n)$, any decomposition into separate data and network effects is underidentified. In the simple additive response $Y=\alpha n+\beta d+\varepsilon$, a linear relation $d=kn$ identifies only $\alpha+k\beta$ from that variation. A factorial intervention varying permissible features independently of participant-network exposure can identify separate causal contrasts under randomization and an adequate interference model. It also needs exclusions ensuring that the data intervention does not directly alter prices or regulatory eligibility through another channel.

Approval and repayment data alone condition on the selected funded pool. Without outcomes or randomized offers for excluded applicants, apparently lower default can reflect exclusion rather than improved predictions for a fixed population. The counterexamples above provide explicit mechanisms for this selection effect. Estimating the full equilibrium frontier additionally needs entry, participation and strategic pricing responses under each policy.

\section*{Status and primary source}
The attempt proves a restricted surplus monotonicity result, explicit failures of inclusion and risk monotonicity, and the exact composition condition for default rates. Endogenous platform scale, equilibrium conduct, privacy technology and empirically identified policy frontiers remain unresolved.

Karen Croxson, Jon Frost, Leonardo Gambacorta and Tommaso Valletti (2022), \emph{Platform-based business models and financial inclusion}, BIS Working Paper 986, Introduction and Conclusion. These parts of the accessible primary PDF were inspected. \url{https://www.bis.org/publications/working-paper-986-platform-based-business-models-and-financial-inclusion.pdf}.

\end{document}
